\chapter{Foundations: from factorials to covers, unconditionally}
\label{ch:foundations}

\begin{chapterabstract}
This chapter is the foundations paper, subsuming and updating the July
2026 manuscript \emph{Cute Goldbach Gaps}. After an elementary factorial
warm-up, we prove unconditionally that $\g$ is unbounded: a primorial
obstruction combined with Dirichlet's theorem produces, for every prime
$P$, an even $\N$ with an explicit partition and
$\g(\N)=\mathrm{nextprime}(P)$. We then formulate the construction as a
congruence-cover problem on the prime offsets, in complete and partial
forms, prove the realisation theorem that upgrades any cover to an
actual Goldbach number, and transfer the
Ford--Green--Konyagin--Maynard--Tao long-gap covering into an
unconditional lower bound on the maximal order of $\g$. The chapter
closes with the first-generation records this machinery produced --- a
$1020$-digit complete-cover baseline, a $237$-digit integer with
$\g=109\,621$, and a $199$-digit integer with $\g=105\,667$ --- which are
the incumbents every later chapter improves.
\end{chapterabstract}

\section{Factorial neighbourhoods and Wilson obstructions}\label{sec:factorial}

\begin{lemma}[Factorial composites]\label{lem:fact}
For $n\ge2$ and $2\le k\le n$, the number $n!+k$ is composite. For
$n\ge4$ and $2\le k\le n$, the number $n!-k$ is also composite.
\end{lemma}
\begin{proof}
In both cases $k\mid n!$, so $k\mid(n!\pm k)$. For the plus side
$n!+k>k$, so $n!+k$ is a proper multiple of $k$. For the minus side,
$n!-k=k\bigl(n!/k-1\bigr)$ with $n!/k-1\ge2$ exactly when $n!\ge3k$,
which holds for $n\ge4$; it fails at $n=3,\,k=3$, where $3!-3=3$ is
prime.
\end{proof}

\begin{lemma}[Wilson obstructions]\label{lem:wilson}
Let $n\ge2$. \emph{(A)} If $n+1$ is prime, then $(n+1)\mid n!+1$.
\emph{(B)} If $n+2$ is prime, then $(n+2)\mid n!-1$.
\end{lemma}
\noindent Recall \emph{Wilson's theorem}: a number $p>1$ is prime if and
only if $(p-1)!\equiv-1\pmod p$, i.e.\ the product $1\cdot2\cdots(p-1)$
is one less than a multiple of $p$. We use only the forward direction.
\begin{proof}
For (A) put $p=n+1$; Wilson's theorem gives $n!=(p-1)!\equiv-1\pmod p$.
For (B) put $p=n+2$; then $(p-1)!=(n+1)\,n!$ and $n+1\equiv-1\pmod p$,
so $-n!\equiv-1$, i.e.\ $n!\equiv1\pmod p$.
\end{proof}

\begin{corollary}\label{cor:facprime}
If $n!+1>n+1$ is prime, then $n+1$ is composite. In particular the
factorial primes $n!+1$ with $n\ge3$ (for example $n=3,11,27,37,41,\dots$)
occur only when $n+1$ is composite.
\end{corollary}

\subsection{The worked example \texorpdfstring{$\N=27!+28$}{N=27!+28}}

Take $\N=27!+28$, even with $29$ digits. Since $27!\equiv0\pmod r$ for
every prime $r\le27$,
\begin{equation}\label{eq:shift}
\N-q\;=\;27!+(28-q)\;\equiv\;28-q\pmod r\qquad(r\le27\ \text{prime}),
\end{equation}
so $\N-q$ is killed by a small prime $r$ exactly when $r\mid(28-q)$: a
shifted sieve that eliminates the class $q\equiv28\pmod r$.

\begin{proposition}\label{prop:71}
For $\N=27!+28$ every prime $q$ with $2\le q\le67$ has $\N-q$ composite,
while $\N-71$ is prime; hence $\g(\N)=71$.
\end{proposition}
\begin{proof}[Proof by explicit witnesses]
Table~\ref{tab:witness} gives, for each prime $q\le71$, a factor of
$\N-q$. For $q\le23$ the witness divides $28-q$ via~\eqref{eq:shift}.
Two cases below the winner are notable: $q=29$ is killed by $29$ because
$29=27+2$ is prime, so Lemma~\ref{lem:wilson}(B) gives
$29\mid27!-1=\N-29$; and $q=59$ survives the sieve (as $59-28=31$ is
$27$-rough) yet $41\mid27!-31=\N-59$. Finally $71-28=43$ is $27$-rough
and $27!-43$ is prime.
\end{proof}

\begin{table}[h]
\centering\small
\caption{Why every small prime fails as a summand of $\N=27!+28$.}\label{tab:witness}
\begin{tabular}{rlc@{\qquad}rlc}
\toprule
$q$ & $\N-q$ & factor & $q$ & $\N-q$ & factor\\
\midrule
$2$&$27!+26$&$2$ & $31$&$27!-3$&$3$\\
$3$&$27!+25$&$5$ & $37$&$27!-9$&$3$\\
$5$&$27!+23$&$23$ & $41$&$27!-13$&$13$\\
$7$&$27!+21$&$3$ & $43$&$27!-15$&$3$\\
$11$&$27!+17$&$17$ & $47$&$27!-19$&$19$\\
$13$&$27!+15$&$3$ & $53$&$27!-25$&$5$\\
$17$&$27!+11$&$11$ & $59$&$27!-31$&$\mathbf{41}$ (generic)\\
$19$&$27!+9$&$3$ & $61$&$27!-33$&$3$\\
$23$&$27!+5$&$5$ & $67$&$27!-39$&$3$\\
$29$&$27!-1$&$\mathbf{29}$ (Wilson) & $71$&$27!-43$&\textbf{prime}\\
\bottomrule
\end{tabular}
\end{table}

By~\eqref{eq:shift} a prime $q$ can even be \emph{tested} only if $q-28$
is $27$-rough, so the viable candidates $29,59,71,89,\dots$ are sparse,
and $71$ is merely the first with a prime complement. Sharper still:
$27!+1$ \emph{is} prime (Corollary~\ref{cor:facprime}), the natural
bridge $q=27$, but $27$ is composite, so one falls through to $71$. For
indices that are simultaneously prime and factorial-prime indices
($n=3,11,37,41$) the bridge works and $\g(n!+(n+1))=n$. Heuristically,
the first viable offsets in this factorial family lie on the scale of
$n$, corresponding to roughly $\log\N/\log\log\N$; the rest of the
chapter produces much larger values.

\section{Primorial deserts, made unconditional}\label{sec:primorial}

The factorial wastes each small prime on a residue class. Taking $\N$ to
be a multiple of a primorial wastes nothing.

\begin{lemma}[Primorial core]\label{lem:primcore}
Let $P$ be prime, let $m_P=\prod_{p\le P}p$ be the \emph{primorial} ---
the product of every prime up to $P$ (for instance
$m_{11}=2\cdot3\cdot5\cdot7\cdot11=2310$) --- and let $\N$ be a multiple
of $m_P$ with $\N>2P$. Then $\N-q$ is composite for every prime
$q\le P$; i.e.\ no prime $\le P$ is a Goldbach summand of $\N$.
\end{lemma}
\begin{proof}
For a prime $q\le P$ we have $q\mid m_P\mid\N$ (as $q$ is one of the
factors of $m_P$), hence $q\mid\N-q$, and $\N-q>P\ge q$, so $\N-q$ is a
proper multiple of $q$ and therefore composite.
\end{proof}

Lemma~\ref{lem:primcore} alone gives only claim~(a) of the standard of
evidence. To reach unboundedness with an \emph{actual} Goldbach number
we add Dirichlet's theorem.

\begin{theorem}[Unconditional unboundedness]\label{thm:dirichlet}
For every prime $P$ put $s=\mathrm{nextprime}(P)$. There are infinitely
many primes $t\equiv-s\pmod{m_P}$, and for any such $t$ the even integer
$\N=s+t$ satisfies $m_P\mid\N$, has the explicit Goldbach partition
$\N=s+t$, and has
\[
\g(\N)\;=\;\mathrm{nextprime}(P)\;>\;P.
\]
Hence $\sup_\N\g(\N)=\infty$ unconditionally. Moreover the least such
$t$ satisfies $\log\N\ll P$ by Linnik's theorem~\cite{Linnik}, so
$\g(\N)\gg\log \N$ along this family.
\end{theorem}
\begin{proof}
As $\gcd(s,m_P)=1$ (because $s>P$), Dirichlet~\cite{Dirichlet} gives
infinitely many primes $t\equiv-s\pmod{m_P}$. Then
$\N=s+t\equiv0\pmod{m_P}$ and $\N$ is even. By Lemma~\ref{lem:primcore}
no prime $\le P$ is a summand; the only primes below
$s=\mathrm{nextprime}(P)$ are those $\le P$, so the least summand is at
least $s$. But $\N-s=t$ is prime, so $s$ is a summand; hence
$\g(\N)=s$. Linnik's theorem bounds the least such prime by
$t\ll m_P^{L}$ for an absolute constant $L$, and
$\log m_P=\theta(P)\sim P$, so $\log\N\ll P\sim s=\g(\N)$.
\end{proof}

\begin{remark}[Concrete instances]\label{rem:concrete}
For $P=13$ one has $m_{13}=30030$, $s=17$, and the least prime
$t\equiv-17\pmod{30030}$ is $t=30013$, giving
$\N=30030=2\cdot3\cdot5\cdot7\cdot11\cdot13$ with partition
$30030=17+30013$ and $\g(\N)=17$ --- all primality deterministic. A
larger fully proven example is the primorial
$m_{43}=13082761331670030$ itself: its minimal partition is
$m_{43}=89+13082761331669941$ (both proven prime, being $<2^{64}$), so
$\g(m_{43})=89$ unconditionally.
\end{remark}

\begin{remark}[The roughness bonus]\label{rem:rough}
For the bare choice $\N=m_P$, a prime $q>P$ makes $\N-q$ coprime to all
primes $\le P$ --- the complement is \emph{$P$-rough} --- which improves
its chances of being prime, since the small primes that catch most
composites are already excluded. By \emph{Mertens' theorem} (which says
$\prod_{p\le P}(1-1/p)\sim e^{-\gamma}/\log P$: the product of $1-1/p$
over primes $p\le P$ shrinks like $1/\log P$), this sieving effect
multiplies the heuristic primality probability by
$\prod_{p\le P}(1-1/p)^{-1}\sim e^\gamma\log P$. The expected least
surviving summand then lands near $P(1+e^{-\gamma})\approx1.56\,P$.
Either way the primorial pays a modulus $m_P$ with
$\log m_P=\theta(P)\sim P$ --- exponentially large in the target --- and
the next section removes this waste. The roughness bonus itself,
however, is not waste but the engine of everything that follows: it
reappears in Chapter~\ref{ch:costmodel} as the ``boost'' factor of the
cost model.
\end{remark}

\section{Prime-offset congruence covers}\label{sec:covering}

\subsection{The construction}

A primorial spends one prime per summand; a prime-offset congruence
cover reuses each prime on a whole residue class. The one structural
constraint is parity: a sum of two odd primes is even, so $\N$ is kept
even and the prime $2$ kills only $q=2$; every \emph{odd} prime below
the target must be covered by odd sieving primes.

The strategy itself is classical. Covering systems of congruences were
introduced by Erd\H{o}s precisely to obstruct an additive
representation: his construction of an infinite arithmetic progression
of odd integers not of the form $2^k+p$~\cite{Erdos1950} is the direct
ancestor of this section (and of the Sierpi\'nski--Riesel tradition that
grew from it). What changes here is only the target: the offsets to be
blocked are the primes below $Q$ rather than the powers of two.

\begin{proposition}[Covering certificate]\label{prop:cover}
Let $Q\ge3$ and let $\{(r,a_r):r\in\mathcal R\}$ be a system of odd
primes $\mathcal R$ such that every odd prime $q<Q$ has
$q\equiv a_r\pmod r$ for some $r\in\mathcal R$. If $\N$ is even,
$\N\equiv a_r\pmod r$ for all $r\in\mathcal R$, and
\[
\N\;>\;Q+\max_{r\in\mathcal R}r,
\]
then $\N-q$ is composite for every prime $q<Q$; i.e.\ no prime $<Q$ is a
Goldbach summand of $\N$.
\end{proposition}
\begin{proof}
For $q=2$, $\N-2$ is even and $>2$. For an odd prime $q<Q$, pick
$r\in\mathcal R$ with $q\equiv a_r\equiv\N\pmod r$; then $r\mid\N-q$,
and $\N-q>(Q+\max_r r)-Q\ge r$, so $r\mid\N-q$ properly and $\N-q$ is
composite.
\end{proof}

\begin{remark}
The size hypothesis only has to make every divisor $r\mid\N-q$ proper,
for which $\N>Q+\max_r r$ suffices; the much stronger $\N>Q\prod_r r$ is
unnecessary. A common solution $\N$ exists modulo $M=2\prod_{r}r$ by the
Chinese Remainder Theorem --- the classical fact that congruences to
pairwise coprime moduli can be satisfied simultaneously, and that the
solutions form exactly one residue class modulo the product --- and any
large enough even representative qualifies.
\end{remark}

We build coverings greedily and deterministically.

\medskip
\noindent\fbox{\parbox{0.97\textwidth}{%
\textbf{Algorithm (greedy covering of the odd primes below $Q$).}
\begin{enumerate}\setlength\itemsep{1pt}
\item Let $T=\{\text{odd primes }<Q\}$ and require $\N\equiv0\pmod2$.
\item For $r=3,5,7,11,\dots$ in increasing order, while
$T\ne\varnothing$: let $a_r$ be a residue mod $r$ containing the most
elements of $T$ (smallest such residue on ties), require
$\N\equiv a_r\pmod r$, and delete from $T$ every $q\equiv a_r\pmod r$.
\item Recover $\N$ from the congruences by the Chinese Remainder
Theorem.
\end{enumerate}}}
\medskip

The resulting modulus is $M=2\prod_{r\in\mathcal R}r$ --- twice the
product of the sieving primes actually used. For the computed targets it
is far smaller than the primorial $\prod_{p<Q}p$, because each sieving
prime is reused across an entire residue class rather than assigned to
one offset. (Chapter~\ref{ch:records} replaces this naive greedy pass
with weighted set cover plus annealing; Chapter~\ref{ch:closure} asks,
and largely answers, how far from optimal such covers are.)

\subsection{Partial covers and progression search}\label{sec:partial}

A complete cover is convenient to verify, but it is not the cheapest way
to make $\N$ small. Near the end of the greedy algorithm, a new sieving
prime may remove only one or two remaining offsets while still
multiplying the modulus by $r$. That is an expensive way to buy a single
compositeness proof.

The alternative is to stop early. Let $U$ be the odd prime offsets below
$Q$ not covered by the chosen congruences --- the \emph{residual set}.
CRT gives one residue $\N_0$ modulo
\[
M=2\prod_{r\in\mathcal R}r,
\]
so every candidate has the form
\begin{equation}\label{eq:progression}
\N=\N_0+kM,\qquad k=0,1,2,\dots
\end{equation}
The fixed congruences dispose of most offsets for every $k$. For each
remaining $q\in U$, we test $\N-q$ directly. A failed probable-prime
test is a rigorous compositeness result: pseudoprimes can create false
positives, but they cannot turn a prime into a reported composite. Thus
a fast one-sided test is safe as a search filter, while the final
surviving complement still needs a primality certificate.

\begin{proposition}[Partial-cover certificate]\label{prop:partial}
Let $U$ be the odd primes $q<Q$ not covered by the congruences of
Proposition~\ref{prop:cover}. Let $\N$ be a sufficiently large even
solution of those congruences. If $\N-q$ is composite for every
$q\in U$, then no prime $q<Q$ is a Goldbach summand of $\N$.
\end{proposition}
\begin{proof}
The congruence witnesses prove compositeness for every covered odd
prime; parity handles $q=2$; and the hypothesis handles the residual
set $U$.
\end{proof}

Most values of $k$ fail the hypothesis --- some residual complement
turns out to be prime, which is not a disaster but a small Goldbach
partition, i.e.\ a losing lottery ticket --- and the construction
becomes a \emph{search} along the progression \eqref{eq:progression}.
What that search costs, and how predictably, is the subject of
Chapter~\ref{ch:costmodel}; it is the single question the rest of the
thesis orbits.

\subsection{Realising a covering unconditionally}

Proposition~\ref{prop:cover} still only gives claim~(a). The same
Dirichlet idea as in \S\ref{sec:primorial}, applied twice, upgrades any
covering to an actual Goldbach number.

\begin{theorem}[Unconditional realisation]\label{thm:coverdirichlet}
Let $Q\ge3$ and let $\{(r,a_r):r\in\mathcal R\}$ cover the odd primes
below $Q$, where every $r\in\mathcal R$ is an odd prime. Then there are
infinitely many even integers $\N>Q+\max_{r\in\mathcal R}r$ admitting an
explicit Goldbach partition $\N=s+t$ with $\N\equiv a_r\pmod r$ for all
$r\in\mathcal R$. Every such $\N$ has $\g(\N)\ge Q$.
\end{theorem}
\begin{proof}
For each $r\in\mathcal R$ choose
$b_r\in\{1,\dots,r-1\}\setminus\{a_r\}$, and also require
$s\equiv1\pmod2$. The CRT gives a reduced residue class $b$ modulo
$2\prod_r r$. Dirichlet gives infinitely many primes in that class; call
one of them $s$. If $s<Q$, then $s$ is an odd covered prime, so
$s\equiv a_r\pmod r$ for some $r$, contradicting
$s\equiv b_r\not\equiv a_r\pmod r$. Hence $s\ge Q$.

Next require
\[
t\equiv a_r-s\pmod r\quad(r\in\mathcal R),\qquad t\equiv1\pmod2.
\]
Every residue $a_r-s$ is nonzero modulo $r$, so this is a reduced
residue class modulo $2\prod_r r$. Dirichlet gives infinitely many prime
values of $t$; keep those large enough that $s+t>Q+\max_r r$, and put
$\N=s+t$. Then $\N$ is even, $\N\equiv a_r\pmod r$, and $\N=s+t$ is an
explicit Goldbach partition. Proposition~\ref{prop:cover} excludes every
prime below $Q$, so $\g(\N)\ge Q$.
\end{proof}

\begin{remark}[Quantitative form for initial-segment covers]\label{rem:size}
The covering used in Theorem~\ref{thm:lower} below has one modulus for
every prime $p\le x$. In that setting the two residue classes above are
reduced modulo $M=2\prod_{p\le x}p$. Linnik's theorem~\cite{Linnik}
supplies primes $s,t\ll M^L$ for an absolute constant $L$. Their
residues are nonzero modulo every prime at most $x$, and the covering
property forces the odd prime $s$ to satisfy $s\ge Q$. Consequently the
resulting integer satisfies
\[
\log\N\ll\log M=\theta(x)\ll x.
\]
This quantitative statement is used only for those initial-segment
covers; no such bound is asserted for an arbitrary sparse finite cover.
\end{remark}

\section{Transfer from large prime gaps}\label{sec:gaptransfer}

\begin{proposition}\label{prop:transfer}
If $[\N-Q,\N-2]$ contains no prime, then no prime $q\le Q$ is a Goldbach
summand of $\N$. Consequently, if $\N$ admits a Goldbach representation,
then $\g(\N)>Q$.
\end{proposition}
\begin{proof}
Every prime $q\in[2,Q]$ has $\N-q\in[\N-Q,\N-2]$, hence $\N-q$ is
composite.
\end{proof}

So a prime gap of length $\ge Q$ just below $\N$ forces $\g(\N)>Q$:
\emph{prime gaps transfer to least Goldbach summands}. The transfer is,
however, strictly one-directional. Forcing $\g(\N)>Q$ requires $\N-q$
composite only for \emph{prime} offsets $q<Q$, not for every integer in
an interval; the covering of Proposition~\ref{prop:cover} covers just
the primes below $Q$, a far sparser target. Thus ``large least Goldbach
summand'' is a relaxation of ``large prime gap,'' and one should not
expect the two problems to coincide. Chapter~\ref{ch:gaps} measures the
separation on certified examples; here we exploit the easy direction.

Writing $G(X)$ for the largest gap between consecutive primes below
$X$, the work of Ford--Green--Konyagin--Tao and of
Maynard~\cite{FGKT,Maynard}, sharpened in~\cite{FGKMT}, gives
\[
G(X)\;\gg\;\frac{\log X\,\log_2 X\,\log_4 X}{\log_3 X},
\]
improving Rankin~\cite{Rankin}. In words: there are stretches of
consecutive composite integers below $X$ of length at least a constant
times $\log X\cdot\log_2 X\cdot\log_4 X/\log_3 X$. The dominant factor
is $\log X$ (the average prime gap near $X$); the three iterated-log
factors are a slowly growing extra boost, and since each $\log_k$ grows
glacially this is only a little more than $\log X$. Feeding the covering
that underlies these bounds into Theorem~\ref{thm:coverdirichlet} (which
costs only a constant factor in $\log\N$, by Linnik) makes the transfer
unconditional:

\begin{theorem}\label{thm:lower}
There are infinitely many even integers $\N$, each with an explicit
Goldbach representation, such that
\[
\g(\N)\;\gg\;\frac{\log \N\,\log_2 \N\,\log_4 \N}{\log_3 \N}.
\]
\end{theorem}

\begin{proof}
The Erd\H{o}s--Rankin method, in the sharpened form of
Ford--Green--Konyagin--Maynard--Tao, supplies the covering. Precisely,
\cite{FGKMT} shows (in the standard covering formulation of the method,
the quantity $Y(x)$ of \cite{FGKT,FGKMT}) that there is an absolute
constant $c>0$ such that for all large $x$ one may choose residue
classes $a_p\pmod p$, one for each prime $p\le x$, whose union contains
every integer in $[1,y]$, where
\[
y=\Bigl\lfloor c\,\frac{x\,\log x\,\log_3 x}{\log_2 x}\Bigr\rfloor;
\]
the prime-gap bound is deduced from exactly this covering by the Chinese
Remainder Theorem.

\emph{Normalisation.} If $a_2$ is odd, replace every class $a_p$ by
$a_p-1$; the shifted system covers $[0,y-1]\supseteq[1,y-1]$. So, at the
cost of replacing $y$ by $y-1$, we may assume $2\mid a_2$. Then no odd
integer in $[1,y-1]$ is covered by the class modulo $2$, so every odd
integer $m\le y-1$ --- in particular every odd prime $q<Q:=y-1$ ---
satisfies $q\equiv a_r\pmod r$ for some \emph{odd} prime $r\le x$. The
pairs $(r,a_r)$ over the odd primes $r\le x$ therefore form a covering
of the odd primes below $Q$ in the sense of
Proposition~\ref{prop:cover}: it covers far more (all odd integers below
$Q$), but only the primes are needed, and self-covering ($r=q$,
$a_r=0$) is harmless, exactly as in the computed systems of
Chapter~\ref{ch:records}.

\emph{Realisation.} Theorem~\ref{thm:coverdirichlet} applied to this
covering yields infinitely many even $\N$ with an explicit partition and
$\g(\N)\ge Q$; by Remark~\ref{rem:size} we may take
\[
\log\N\;\ll\;\log2+\theta(x)\;\ll\;x.
\]

\emph{Conclusion.} Let $C$ be the absolute constant with
$\log\N\le Cx$. For large $u$ the function
$u\mapsto u\log u\,\log_3u/\log_2u$ is increasing, and each iterated
logarithm of $\log\N/C$ is $(1+o(1))$ times the corresponding iterated
logarithm of $\log\N$; hence
\begin{align*}
\g(\N)\;\ge\;Q\;&\gg\;\frac{x\,\log x\,\log_3x}{\log_2x}
\;\ge\;\frac{\tfrac{\log\N}{C}\,\log\bigl(\tfrac{\log\N}{C}\bigr)\,\log_3\bigl(\tfrac{\log\N}{C}\bigr)}{\log_2\bigl(\tfrac{\log\N}{C}\bigr)}\\[2pt]
&\gg\;\frac{\log\N\,\log_2\N\,\log_4\N}{\log_3\N}.
\end{align*}
For each sufficiently large $x$, select one realisation $\N_x$
satisfying this Linnik bound. Since $\g(\N_x)\ge Q(x)\to\infty$, the
selected integers are unbounded and therefore contain infinitely many
distinct values. This proves the claim.
\end{proof}

Together with the Granville--van de Lune--te Riele
conjecture~\eqref{eq:GLR}, this brackets the maximal order of $\g$
between a logarithmic-type lower bound and a $(\log\N)^2$-type
conjectural ceiling:
\[
\underbrace{\frac{\log \N\,\log_2 \N\,\log_4 \N}{\log_3 \N}}_{\text{rigorous lower-bound scale}}
\;\ll\;\max_{\N'\le \N}\g(\N')
\;\lesssim\;
\underbrace{(\log \N)^2\log\log \N}_{\text{conjectural ceiling}}.
\]
The explicit constructions of this thesis do not settle that gap; their
value is that they turn the problem into auditable objects --- a
congruence list, a CRT progression, a finite residual set, and a
primality certificate --- and, in Chapters~\ref{ch:limits}
and~\ref{ch:closure}, into measured statements about which parts of the
gap are closable at all. We note one reappearance of this asymptotic
machinery: the FGKMT covering enters again in
Chapter~\ref{ch:closure} not as a lower bound but as the one known
\emph{rounding technology} that might beat greedy covers, which is this
thesis's central open problem.

\section{The first-generation records}\label{sec:firstgen}

The machinery above produced the first constructive records in this
problem family, which the rest of the thesis treats as its baseline.
All three claims (a)/(b)/(c) were discharged to the standard of
Chapter~\ref{ch:intro}, and all were re-verified inside the thesis
pipeline before any successor record was accepted.

\paragraph{The complete-cover baseline.} A complete greedy cover of the
odd primes below $Q=10^5$ uses $356$ congruences and forces a modulus of
$1020$ digits; its least even representative has $\g(\N)=100\,747$.
Every excluded offset has a small factor visible directly from the
congruence system, so verification is little more than modular
arithmetic --- a \emph{proof-oriented build} that overpays roughly
sevenfold in digits for certainty it does not need. (A complete cover at
$Q=10^8$ was also computed as a scaling probe: $327\,455$ digits, with
no partition exhibited.)

\paragraph{The hybrid ablation.} Stopping the cover early and searching
the progression \eqref{eq:progression} shrinks the integer dramatically.
The engineering path, at fixed target $Q=10^5$:

\begin{table}[h]
\centering\small
\caption{First-generation ablation at target $Q=10^5$. ``Congr.''
includes parity; ``residual'' counts prime offsets below $10^5$ left
for direct testing; the last column is the first prime offset whose
complement is (certified) prime.}\label{tab:ablation}
\begin{tabular}{lrrrrrr}
\toprule
Method & Congr. & Digits $M$ & Residual & $k$ & Digits $\N$ & $\g(\N)$\\
\midrule
Complete greedy cover & 356 & 1020 & 0 & 0 & 1020 & $100\,747$\\
Prefix partial cover & 168 & 416 & 385 & 74 & 418 & $101\,599$\\
Coordinate descent & 110 & 248 & 580 & $1\,751$ & 251 & $100\,981$\\
Sparse weighted cover & 103 & 232 & 610 & $299\,948$ & 237 & $109\,621$\\
\bottomrule
\end{tabular}
\end{table}

\paragraph{The $237$-digit height record.} The final row is the
first-generation height incumbent: a $237$-digit even integer
$\N_{237}=\N_0+299\,948\,M$ over a partial cover of $102$ odd
congruences plus parity ($M$ of $232$ digits, largest modulus $1217$),
with
\[
\g(\N_{237})=109\,621.
\]
Every prime offset below $109\,621$ is excluded --- $9\,691$ by explicit
congruence divisors and $731$ by strong base-2 compositeness witnesses
--- and the complement $\N_{237}-109\,621$ is ECPP-certified prime.

\paragraph{The $199$-digit dual incumbent.} Re-optimising the cover for
the strict budget $\N<10^{199}$ gives a $199$-digit integer
$\N_{199}=\N_0+876\,554\,M$ ($89$ odd congruences plus parity, $M$ of
$193$ digits) with
\[
\g(\N_{199})=105\,667,
\]
proved by $9\,335$ parity-or-congruence witnesses, $744$ strong base-2
witnesses, and a $34$-step elliptic-curve certificate for the
complement. It entered the thesis as the incumbent for \emph{both}
bounded games: $T(100\,000)\le\N_{199}$ and $R(10^{199})\ge105\,667$.

\medskip
These records made the covering paradigm concrete, and they framed the
questions the thesis then had to answer. Why $237$ digits and not $150$?
Is the greedy cover any good? What does a progression search actually
cost, and can that cost be predicted rather than experienced? The next
chapter builds the measuring instruments; Chapter~\ref{ch:records}
spends them.
